\documentclass[11pt]{amsproc}

\usepackage[a4paper,margin=2.6cm]{geometry}
\usepackage{amsmath,amssymb,amsthm}
\usepackage[T1]{fontenc}
\usepackage[utf8]{inputenc}

\newcommand{\C}{\mathbb{C}}
\newcommand{\gen}[1]{\langle #1\rangle}
\newcommand{\col}{\operatorname{col}}

\begin{document}

\section*{Problema}

Sea $V=\mathbb{R}^n$ y sea $ G$ una familia finita, con al menos dos
elementos, de transformaciones lineales invertibles de $V$, cerrada bajo
composición.
Sean $v_1,\dots,v_n\in V$ vectores linealmente independientes. Para cada
$i=1,\dots,n$, definimos
\[
    w_i=\frac1{| G|}\sum_{T\in G}T(v_i).
\]
Calcule el volumen del simplex con vértices
$0,w_1,\dots,w_n$.

\section*{Solución}
Como $ G$ es finita, está formada por transformaciones invertibles y
es cerrada bajo composición, $ G$ es un grupo.
Sea
\[
    S=\sum_{T\in G}T.
\]
Para cualquier $U\in G$ tenemos
\[
    US=\sum_{T\in G}UT
      =\sum_{T\in G}T
      =S.
\]

Como $ G$ tiene al menos dos elementos, podemos elegir $U\in G$ distinto de la
identidad. De $US=S$ se sigue que $(U-I)S=0$. Si $S$ es inversible, entonces
$U=I$, una contradicción. Luego $S$ no es inversible y $\det S=0$.  Sea $A$ la
matriz cuyas columnas son $v_1,\dots,v_n$. La matriz cuyas columnas son
$w_1,\dots,w_n$ es
\[
    \frac1{| G|}SA.
\]
Entonces, 
\[
\det(w_1,\dots,w_n)
    =\frac{1}{| G|^n}\det(S)\det(A)
    =0.
\]
Por lo tanto, el simplex con vértices $0,w_1,\dots,w_n$ tiene volumen cero. 

\section*{Problema (una variante del anterior)}
Sea $ G$ una familia finita de $m$ transformaciones lineales
invertibles de $\mathbb R^n$, cerrada bajo composición. Sean
$v_1,\dots,v_n$ vectores linealmente independientes y, para cada $i$, defina
\[
    w_i=v_i+\sum_{T\in G}T(v_i).
\]
Demuestre que el cociente entre los volúmenes de los simplex
\[
    (0,w_1,\dots,w_n)
    \qquad\text{y}\qquad
    (0,v_1,\dots,v_n)
\]
es una potencia de $m+1$.

\section*{Solución}

Como $ G$ es finita, está formada por transformaciones invertibles y
es cerrada bajo composición, $ G$ es un grupo.
Consideremos el operador
\[
    S=\sum_{T\in G}T.
\]
Para todo $U\in G$ se tiene
$US=S$, 
pues la multiplicación por $U$ simplemente permuta los elementos de
$ G$. En particular, la imagen de $S$ está contenida en el subespacio
\[
    W=\{v\in\mathbb R^n:T(v)=v\text{ para todo }T\in G\}.
\]
Por otra parte, si $v\in W$, entonces
$S(v)=mv$. 
Más aún,
\[
    S^2
    =\sum_{T,U\in G}TU
    =m\sum_{R\in G}R
    =mS,
\]
pues cada elemento $R\in G$ aparece exactamente $m$ veces en la
doble suma.
Por lo tanto, el operador
$P=\frac1mS$ 
satisface $P^2=P$. Es decir, $P$ es una proyección, y su imagen es
precisamente $W$.
Sea $r=\dim W$. Los autovalores de $P$ son entonces
\[
    \underbrace{1,\dots,1}_{r\text{ veces}},
    \underbrace{0,\dots,0}_{n-r\text{ veces}}.
\]
En consecuencia, los autovalores de
$I+S=I+mP$ 
son $m+1$, con multiplicidad $r$, y $1$, con multiplicidad $n-r$. Así,
$\det(I+S)=(m+1)^r$. 
Ahora bien,
$w_i=(I+S)v_i$.
Si $A$ es la matriz cuyas columnas son $v_1,\dots,v_n$, la matriz cuyas
columnas son $w_1,\dots,w_n$ es
$(I+S)A$.
Por consiguiente,
\[
    \frac{\operatorname{Vol}(0,w_1,\dots,w_n)}
         {\operatorname{Vol}(0,v_1,\dots,v_n)}
    =
    |\det(I+S)|
    =(m+1)^r.
\]

\end{document}
